\documentclass[12pt]{article}
\usepackage[utf8]{inputenc}
\usepackage[spanish]{babel}
\usepackage[a4paper, margin=2cm]{geometry}
\usepackage{tikz}
\usetikzlibrary{babel}
\usepackage[most]{tcolorbox}
\usepackage{amsmath}

\renewcommand{\familydefault}{\sfdefault}

% Definición de colores
\definecolor{medblue}{RGB}{43, 108, 176}
\definecolor{medteal}{RGB}{49, 151, 149}
\definecolor{medlight}{RGB}{235, 248, 255}
\definecolor{meddark}{RGB}{26, 32, 44}
\definecolor{heartred}{RGB}{220, 38, 38}    % Rojo para el corazón
\definecolor{ecggreen}{RGB}{16, 185, 129}   % Verde brillante para el ECG
\definecolor{powerorange}{RGB}{245, 158, 11} % Naranja para la energía

\begin{document}
\pagestyle{empty}

\begin{tcolorbox}[
    enhanced, 
    colback=medlight!30, 
    colframe=medblue, 
    title={\Large \textbf{Potencia Mecánica: Gasto Cardíaco}}, 
    halign title=center, 
    fonttitle=\bfseries, 
    arc=4mm, 
    boxrule=1.5pt,
    shadow={2mm}{-2mm}{0mm}{black!30!white}
]

    \vspace{0.2cm}
    \begin{tcolorbox}[colback=white, colframe=medteal, arc=2mm, boxrule=1.2pt, title=\textbf{Caso Clínico / Problema}]
        El ventrículo izquierdo de un corazón humano en reposo realiza aproximadamente \textbf{1 J} de trabajo mecánico por cada latido. Si la frecuencia cardíaca del paciente es de \textbf{72 latidos por minuto}, ¿cuál es la potencia media desarrollada por el ventrículo en Watts?
    \end{tcolorbox}

    \vspace{0.4cm}
    
    \begin{center}
    \begin{tikzpicture}[scale=0.95]
        
        % --- Corazón (Izquierda) ---
        \begin{scope}[shift={(-4, 0.5)}, scale=0.55]
            \filldraw[heartred, draw=meddark, line width=1.5pt] 
                (0,0.5) .. controls (-1.5,2.2) and (-3.5,-0.5) .. 
                (0,-2.5) .. controls (3.5,-0.5) and (1.5,2.2) .. (0,0.5);
            \node[font=\bfseries, text=white] at (0,-0.6) {1 J};
            \node[font=\bfseries, text=meddark, align=center] at (0,-3.8) {Trabajo por\\latido};
        \end{scope}

        % Flecha 1
        \draw[->, ultra thick, medteal] (-2.2, 0) -- (-1.6, 0);

        % --- ECG / Frecuencia (Centro) ---
        \begin{scope}[shift={(0.8, 0)}]
            % Pantalla del monitor
            \filldraw[meddark!95!black, draw=meddark, thick, rounded corners=3pt] (-2.2, -1.2) rectangle (2.2, 1.2);
            % Cuadrícula sutil
            \draw[step=0.2, gray!40!meddark, very thin] (-2.2,-1.2) grid (2.2,1.2);
            
            % Trazo ECG (Ondas P, QRS, T) repetido dos veces
            \draw[ecggreen, thick] 
                (-2.2, 0) -- (-1.7, 0) -- (-1.5, 0.2) -- (-1.4, -0.1) -- (-1.2, 0.8) -- (-1.0, -0.3) -- (-0.8, 0) -- (-0.4, 0) -- (-0.2, 0.3) -- (0, 0) --
                (0.3, 0) -- (0.5, 0.2) -- (0.6, -0.1) -- (0.8, 0.8) -- (1.0, -0.3) -- (1.2, 0) -- (1.6, 0) -- (1.8, 0.3) -- (2.0, 0) -- (2.2, 0);
                
            \node[font=\bfseries, text=meddark, align=center] at (0, -2.1) {Frecuencia Cardíaca\\72 latidos / 60 s};
        \end{scope}

        % Flecha 2
        \draw[->, ultra thick, medteal] (3.2, 0) -- (3.8, 0);

        % --- Potencia / Resultado (Derecha) ---
        \begin{scope}[shift={(5.5, 0)}]
            \filldraw[powerorange!20, draw=powerorange, thick, rounded corners=5pt] (-1.2, -1.2) rectangle (1.2, 1.2);
            \node[font=\Huge\bfseries, text=powerorange!80!black] at (0, 0.1) {1.2};
            \node[font=\large\bfseries, text=powerorange!80!black] at (0, -0.6) {Watts};
            \node[font=\bfseries, text=meddark, align=center] at (0, -2.1) {Potencia Media\\($P = W / t$)};
        \end{scope}
        
    \end{tikzpicture}
    \end{center}

    \vspace{0.4cm}
    \begin{tcolorbox}[colback=medteal!10, colframe=medteal, arc=2mm, boxrule=1.2pt, title=\textbf{Desarrollo Matemático y Solución}]
        La potencia ($P$) es la rapidez con la que se realiza un trabajo mecánico, expresada como $P = \frac{W}{t}$. Para obtener la potencia en Watts (Joules por segundo), primero calculamos el trabajo total realizado en un minuto (60 segundos):
        \[ W_{\text{total}} = 72 \text{ latidos} \times 1 \text{ J/latido} = 72 \text{ Joules} \]
        \[ P = \frac{72 \text{ J}}{60 \text{ s}} = 1.2 \text{ J/s} = \mathbf{1.2 \text{ Watts}} \]
        \textbf{Resultado:} El ventrículo izquierdo desarrolla una potencia media de \textbf{1.2 W} en reposo.
    \end{tcolorbox}
    
\end{tcolorbox}

\end{document}